% Requires amsmath, amssymb, amsthm, hyperref and theorem/proof environments.
\section{Combining the two newest equational reductions: E667}
\label{sec:e667-combination}

The two incoming packets can be combined. The 692-bit packet in the
ProveIt input at commit
\texttt{ecb4df12d256704d9115d4860f6ccb7b3443ef4a}
shares three witness constructors, optimizes the Boolean hull expression,
and reuses the unary class binder as the successor of a selector.
The subsequent incoming 236-state archive, added at commit
\texttt{0b5fec6980cc3a5ba7dd47b404f0cbbe9a451c02}, independently derives
the truth equation of E713 from its other equations, obtaining E688.
Applying that \emph{derived-truth argument} to E692 gives the complete
\textbf{667-bit packet E667} below.

This combination requires a fresh algebraic check. E692 changed the
closure polynomial using an identity proved under the old truth axiom;
it would be circular simply to delete that axiom while citing the
earlier optimized-closure equivalence. Instead, we prove directly that
the \emph{optimized} closure still forces truth without assuming it.

\subsection{The exact seven equations}

The unchanged ordered signature is
\[
 d^0,N^2,V^2,P^2,A^2,Q^1,W^2,p_E^2,p_\varepsilon^2.
\]
The first six equations, with all displayed variables universally
quantified, are
\begin{align}
 N(N(x,N(N(y,x),x)),N(y,N(z,x)))&=y,
     \label{eq:e667-sh1}\\
 V(P(x,s),d)&=x,\label{eq:e667-head}\\
 V(P(x,s),Q(j))&=V(s,j),\label{eq:e667-tail}\\
 p_\varepsilon(s,P(a,b))
   &=N(p_\varepsilon(s,a),p_\varepsilon(s,b)),
       \label{eq:e667-nand}\\
 p_\varepsilon(s,A(i,j))&=p_E(V(s,i),V(s,j)),
       \label{eq:e667-atom}\\
 p_\varepsilon(s,Q(a))
   &=N\bigl(p_\varepsilon(P(W(s,a),s),a),
            p_\varepsilon(P(z,s),a)\bigr).
       \label{eq:e667-binder}
\end{align}
The first is the published Sh1 single axiom for Boolean algebra with
$N$ interpreted as NAND. Its Boolean-algebra consequence is inherited
from the primary proof of McCune, Veroff, Fitelson, Harris, Feist and Wos,
\emph{Short Single Axioms for Boolean Algebra}, Journal of Automated
Reasoning 29 (2002), 1--16; explicit proof files are supplied by the
authors at \url{https://www.cs.unm.edu/~veroff/BA/}.
The symbols $0_B,1_B$ below mean the \emph{derived} Boolean bottom and
top of this reduct. They are not extra symbols in the packet, and $d$
is not initially identified with either of them.

The seventh equation uses the following abbreviations solely for print
readability. Let
\[
 h=Q(b),\qquad n_z=P(b,z),\qquad m_z=A(b,z),
\]
\[
 R_b(x,y)=p_\varepsilon(P(y,P(x,b)),V(b,d)),
\]
and define the eight Boolean arguments
\begin{align*}
 r&=R_b(h,n_z),&e&=p_E(m_z,n_z),&f&=p_E(m_z,z),\\
 a_0&=p_E(W(x,d),x),&b_0&=p_E(x,h),&c_0&=R_b(x,y),\\
 u&=p_E(y,h),&v&=R_b(P(b,y),y).&&
\end{align*}
The symbols $a_0,b_0,c_0$ distinguish these Boolean values from the
native variables. Define the pure NAND polynomials
\begin{align*}
 K_0&=N(c_0,N(a_0,N(d,b_0))),\\
 H&=N\bigl(N(N(d,u),K_0),N(u,v)\bigr),\\
 B^-&=N(r,N(d,N(e,N(d,f)))),\\
 G(d,r,e,f,a_0,b_0,c_0,u,v)&=N(B^-,N(d,H)).
\end{align*}
The seventh equation is exactly
\begin{equation}
                         d=G(d,r,e,f,a_0,b_0,c_0,u,v).
                         \label{eq:e667-closure}
\end{equation}
All abbreviations are expanded in \texttt{E667.equations}, in the JSON
term trees, and in the raw bit packet. The ground query is unchanged:
\[
                             N(d,d)=d.
\]

Relative to E692, the only deletion is
\begin{equation}
                            d=N(d,N(d,d)).
                            \tag{T}\label{eq:e667-deletedtruth}
\end{equation}
In a Boolean algebra its right side is $1_B$. The arities, symbol order,
other seven equations, variable capacity and both ground endpoints are
identical to those of E692. In particular the unary selector equation
\eqref{eq:e667-tail} and all three structural witness identifications
are preserved.

\subsection{Why the optimized closure still forces truth}

\begin{lemma}[Cofactor of the optimized closure]
In every Boolean algebra, for arbitrary values of all nine displayed
arguments,
\[
 \neg d\wedge G(d,r,e,f,a_0,b_0,c_0,u,v)=\neg d\wedge r.
 \tag{Cofactor}\label{eq:e667-cofactor}
\]
Consequently $r=1_B$ and $d=G(d,r,e,f,a_0,b_0,c_0,u,v)$ imply $d=1_B$.
\end{lemma}
\begin{proof}
Set the formal Boolean argument $d$ to $0_B$. The right input of the
outermost NAND becomes $N(0_B,H)=1_B$, whatever value the optimized hull
polynomial $H$ has. Inside $B^-$, the outer $N(0_B,\cdot)$ likewise
gives $1_B$, so $B^-=N(r,1_B)=\neg r$. Thus
\[
 G(0_B,r,e,f,a_0,b_0,c_0,u,v)=N(\neg r,1_B)=r.
\]
This computation neither uses nor needs an interpretation of the
optimized hull as the original hull when $d\ne1_B$.

For an all-algebra proof of the cofactor identity, let $q=\neg d$.
The map $t\mapsto t\wedge q$ is a Boolean homomorphism to the principal
interval $[0_B,q]$, using its relative complement and relative top $q$.
It sends $d$ to $0_B$. Apply the preceding polynomial calculation in
that interval to obtain \eqref{eq:e667-cofactor}. This projects only
the already evaluated Boolean arguments; it makes no assertion that
any native constructor or predicate commutes with the projection.
If $r=1_B$ and $d=G$, the identity gives
$0_B=\neg d\wedge d=\neg d\wedge1_B=\neg d$, hence $d=1_B$.
\end{proof}

\begin{lemma}[A ground closure instance has $r=1_B$]
The retained head and NAND-class equations provide a ground instance of
\eqref{eq:e667-closure} whose first matrix argument is Boolean top,
without assuming equation~\eqref{eq:e667-deletedtruth}.
\end{lemma}
\begin{proof}
Use the same syntactic constant-true class as the incoming E688 proof:
\[
 k=P(d,P(d,d)),\qquad b_*=P(k,d).
\]
For any native argument $s$, two applications of
\eqref{eq:e667-nand} yield
\[
 p_\varepsilon(s,k)=N\bigl(p_\varepsilon(s,d),
             N(p_\varepsilon(s,d),p_\varepsilon(s,d))\bigr)=1_B.
\]
The last equality is the Boolean identity $N(t,N(t,t))=1_B$.
It does not assume any particular class value for $d$.
Equation~\eqref{eq:e667-head} gives $V(b_*,d)=k$.

Now substitute $b=b_*$ and $x=y=z=d$ in the seventh equation. Its
first matrix value is
\[
 r=p_\varepsilon\bigl(P(P(b_*,d),P(Q(b_*),b_*)),k\bigr)=1_B.
\]
The structural identifications $\mathsf{Target}=Q$ and
$\mathsf{Next}=P$ merely determine the first argument of
$p_\varepsilon$; the displayed constant-class identity holds for every
such argument. The selector-tail equation is not needed in this truth
derivation. Thus neither constructor sharing nor the unary-selector
change invalidates the witness.
\end{proof}

\begin{theorem}[Exact model equivalence of E667 and E692]
\label{thm:e667-equivalence}
E667 and E692 have exactly the same algebras in their common signature,
and hence exactly the same equational consequences. With the explicitly
inherited E692 calibration, ordinary ZFC proves
\[
 \mathcal E_{667}\nvdash N(d,d)=d
 \quad\Longleftrightarrow\quad
 \mathcal E_{692}\nvdash N(d,d)=d
 \quad\Longleftrightarrow\quad
 \operatorname{Con}(\mathrm{ZFC}).
\]
\end{theorem}
\begin{proof}
Every E692 algebra satisfies E667 by deleting its truth axiom.
Conversely, in an E667 algebra Sh1 supplies the Boolean reduct. The
ground instance in the second lemma and the cofactor lemma force
$d=1_B$, which is precisely the missing truth equation. All remaining
E692 equations are retained literally, so this is an E692 algebra.
The argument includes the degenerate Boolean algebra as well.

Equational soundness and completeness turn equality of model classes
into equality of consequences. Equivalently, the derived equation can
replace every use of the old truth axiom in a finite equational proof.
All formula trees, algebras, polynomial calculations and term models
used here are set-sized ZFC objects. Composition with the existing
E692/E713/raw-EPSC/ZFC comparison gives the displayed conclusion.
The raw-EPSC interpretation and the structural-sharing model theorem
are inherited dependencies, not reproved by the finite codec checks.
\end{proof}

The incoming E688 theorem alone would not justify this combination:
its closure polynomial differs from the one in E692. The direct
cofactor calculation above supplies that missing compatibility check.
No claim is made that the optimized and original closure polynomials
agree for arbitrary $d$; only their shared truth-forcing property is
needed before applying the already established E692 calibration.

\subsection{Complete accounting and executable artifacts}

The retained equation-node counts, including both sides, are
\[
                     14,6,9,12,12,17,80,
\]
for a total of 150 nodes. Nine symbols and four variable slots still
require a four-bit term opcode. The unchanged codec charges an
eight-bit symbol count, three bits for each arity, an eight-bit variable
capacity, one bit before each equation, a zero equation-list terminator,
and both ground endpoint trees. Thus
\[
 \underbrace{8+3\cdot9+8}_{43\text{ header}}
 +\underbrace{7+1+4\cdot150}_{608\text{ equations}}
 +\underbrace{4\cdot4}_{16\text{ endpoints}}
 =\boxed{667\text{ bits}}.
\]
The reduction from E692 is exactly the six-node truth equation plus its
equation marker: $6\cdot4+1=25$ bits. Relative to E688, the additional
21 bits are the earlier signature, optimized-hull and unary-selector
savings. These comparisons count complete packets under one codec.

The raw 667 bit characters, excluding the final newline, have SHA-256
\begin{center}\small\ttfamily
1075854399e1c3428e1c178b51bac47e428b8ef010b59e0312a3c06d409e55b8.
\end{center}
The packed transport has 84 bytes, with five trailing zero padding bits;
\texttt{E667.bitlength} records 667 so that a transport adapter can
discard that padding. These five bits are not part of the logical
packet.

Run \texttt{python3 -O build\_e667.py} from the supplied equational
directory. The checker uses explicit exceptions, so all checks remain
enabled under optimization. It pins the E692 JSON and bit stream,
reconstructs its optimized closure independently, checks the unchanged
head/NAND-class/unary-selector equations, deletes only the truth row,
and roundtrips the new packet with the independently implemented codec
adapted from the incoming E688 builder. It rejects all 667 strict
prefix truncations and three malformed controls. It also checks the
cofactor on all 512 assignments to its nine independent Boolean
arguments and verifies the packed-byte transport.

The output includes the full equation trees, expanded equations,
raw and packed bits, bit length, a concrete ground-substitution recipe
for the truth derivation, and a verification receipt. The finite
Boolean table verifies the polynomial identity; the written principal-
interval argument establishes its arbitrary-Boolean-algebra meaning.
A fully expanded Sh1-only contextual rewrite trace, a Lean/Rocq
endpoint, and an accounted binary ground-rewrite interpreter are not
supplied. E667 is a smaller precisely specified consistency problem;
its 667-bit size does not by itself imply a Turing-machine state bound.
